\documentclass[../../../main.tex]{subfiles}

\begin{document}

\section{Pointers}\label{sec:26sbpfz}

Pointers have its reputation of being hard to understand when it is
just another data type. I was planning to not create a separate
section for pointer and just put it in the previous section alongside
other data types, but I wanted to discuss more about arrays and
strings that we missed earlier and discuss about dynamic memory
allocation. So in this section we are going to not only discuss
about pointers, but also related topics.

So pointers in short is just a variable that stores the address
of another variables. So what does this mean? When we store a
value like \verb|int a = 5|, computers have to store that information
somewhere. Integers like these are usually 4 bytes and written in
binaries. Computers store this information in memory, or better
known as RAM (Random Access Memory). When storing the information,
they don't just store it, but allocates a space, or address, in the
memory for the value. Here is an example.

\begin{center}
\begin{tabular}{|c|c|}
    \hline
    \textbf{Memory Address} & \textbf{Value} \\
    \hline
    0x16bc6ac18 & 5 \\
    \hline
    0x16bc6ac10 & 0.100000 \\
    \hline
    0x16f7a2c0f & A \\
    \hline
\end{tabular}
\end{center}

Pointer is just a variable that stores that memory address of an
object. Here is the snippet that declares and prints pointer values.

\begin{cbox}{26sbpfz-01}{26sbpfz-01}
\begin{lstlisting}[style=cstyle]
#include <stdio.h>

int main() {
    int a = 5;
    double b = 0.1;
    char c = 'A';

    int* pA = &a;
    double* pB = &b;
    char* pC = &c;

    printf("%d, %p\n", a, pA);
    printf("%f, %p\n", b, pB);
    printf("%c, %p\n", c, pC);

    return 0;
}
\end{lstlisting}
\end{cbox}
\begin{outputbox}
\begin{lstlisting}[style=basic]
5, 0x16bc6ac18
0.100000, 0x16bc6ac10
A, 0x16f7a2c0f
\end{lstlisting}
\end{outputbox}

Here I used \verb|int* pA|, but \verb|int * pA| and \verb|int *pA|
does the same thing. Also it is a convention to have ``p'' in front
of a pointer variable!

Being able to directly manipulate how the program manages memory
is great because we can efficiently manage memory with manual
allocation, pass data more efficiently by pointing instead of
having a copy of the variable, and build data structures like trees
and linked lists.

We can access the variable with its memory address, or pointer,
and change it. This gives us freedom on how we use memory in our
program. But this freedom is also what makes it easy for us to shoot
ourselves in the foot. But before getting to much into that,
let's take a look at how we can \textbf{dereferece} a pointer to
change the value that it is pointing at.

\begin{cbox}{26sbpfz-02}{26sbpfz-02}
\begin{lstlisting}[style=cstyle]
#include <stdio.h>

int main() {
    int num = 5;
    printf("Original num: %d\n", num);

    int* pNum = &num;
    *pNum = 10;
    printf("New num: %d\n", num);

    return 0;
}
\end{lstlisting}
\end{cbox}
\begin{outputbox}
\begin{lstlisting}[style=basic]
Original num: 5
New num: 10
\end{lstlisting}
\end{outputbox}

As you can see, I like to see it as negative times negative is
positive. We give one more asterisk to let the compiler know that
we are dereferencing a referenced value. With this, I think it is
now evident why we need \& when using \verb|scanf|! We are basically
passing pointers to indirectly change the value of the declared
variable. Let's take a deeper look at how arrays and strings work
with pointers.

A string is literally an array of character. So if we define
\verb|char str[] =| \verb|"Hello"|, it will create an array like
\verb|{'H', 'e', 'l', 'l', 'o',| \verb|'\0'}| where \verb|\0| is
a null character that terminates the string. For that reason,
when we use pointer with a string, it will work index by index.
Here is what I mean.

\begin{cbox}{26sbpfz-03}{26sbpfz-03}
\begin{lstlisting}[style=cstyle]
#include <stdio.h>

int main() {
    char str[] = "Hello";
    char* ptr = str;

    *ptr = 'A';
    printf("%s\n", str);

    *(ptr + 2) = 'A';
    printf("%s\n", str);

    return 0;
}
\end{lstlisting}
\end{cbox}
\begin{outputbox}
\begin{lstlisting}[style=basic]
Aello
AeAlo
\end{lstlisting}
\end{outputbox}

As you can see, when we initialize the pointer for the string,
the pointer is actually pointing to the address of the first
character. The pointer can't store the address of the string because
the string has many addresses. When we want to access diffferent
element in the string, we can use arithmetic to find the right index.
Also, if you wondered why we don't put ampersand for string when
declaring a pointer and using \verb|scanf|, that is because the
pointer is automatically assigned to the first character.

Now as we can see from \verb|char name[20]| example, we can manually
allocate memory in C. This causes a problem when we don't do it
properly. Sometimes if we don't release memory after we no longer
needed, the program will not only degrade performance but also could
crash due to \textbf{memory leaks}. To avoid it, we need to free
memory after allocating it. Here is an example.

\begin{cbox}{26sbpfz-04}{26sbpfz-04}
\begin{lstlisting}[style=cstyle]
#include <stdio.h>
#include <stdlib.h>

int main() {
    int* num = malloc(5 * sizeof(int));

    for (int i = 0; i < 5; i++) {
        num[i] = i;
    }
    for (int i = 0; i < 5; i++) {
        printf("%d ", num[i]);
    }
    printf("\n");

    free(num);

    return 0;
}
\end{lstlisting}
\end{cbox}
\begin{outputbox}
\begin{lstlisting}[style=basic]
0 1 2 3 4
\end{lstlisting}
\end{outputbox}

Without \verb|free(num);|, we will have memory leaks and cause
degradation of the program. Continuing from the strings, you may have
seen lots of arrows in C programs. Those are actually pointers as
well! Here is an example of using pointers with structs.

\begin{cbox}{26sbpfz-05}{26sbpfz-05}
\begin{lstlisting}[style=cstyle]
#include <stdio.h>

typedef struct {
    int a;
    int b;
    int c;
} array;

int main() {
    array arr = { 1, 2, 3 };
    array* parr = &arr;
    printf("Original Array: [%d, %d, %d]\n", arr.a, arr.b, arr.c);

    (*parr).a = 4;
    parr->b = 5;
    parr->c = 6;
    printf("New Array: [%d, %d, %d]\n", arr.a, arr.b, arr.c);

    return 0;
}
\end{lstlisting}
\end{cbox}
\begin{outputbox}
\begin{lstlisting}[style=basic]
Original Array: [1, 2, 3]
New Array: [4, 5, 6]
\end{lstlisting}
\end{outputbox}

Finally before wrapping up this section, we can also now write and
modify files with C using pointers!

\begin{cbox}{26sbpfz-06}{26sbpfz-06}
\begin{lstlisting}[style=cstyle]
#include <stdio.h>

int main() {
    FILE* fptr = fopen("file.txt", "w"); // r, w, a
    fprintf(fptr, "Hello World\n");
    fclose(fptr);

    return 0;
}
\end{lstlisting}
\end{cbox}

Running this, we should get a file ``file.txt'' in the directory where
the program was ran with the text ``Hello World'' in it. This is it
for pointers! Before we write exciting programs, we will be going
over header guard and macro in the next section.

\end{document}


